(1) 单比特态写为 \(\rho=(I+\boldsymbol s\cdot\boldsymbol\sigma)/2\)，\(|\boldsymbol s|\le1\)。取四面体单位矢量 \(\boldsymbol r_j=(\pm1,\pm1,\pm1)/\sqrt3\)，选符号积为正的四个顶点，使 \(\boldsymbol r_i\cdot\boldsymbol r_j=-1/3\)（\(i\ne j\)）。定义四效果 \(E_j=(I+\boldsymbol r_j\cdot\boldsymbol\sigma)/4\)，证明这是合法且信息完备的POVM。
(2) 从 Born 规则求单次结果概率 \(p_j\)，构造 Bloch 向量的严格线性反演。明确由统计频率直接得到的向量何时可能落在物理 Bloch 球外。
(3) 对 \(N\) 次相互独立的相同准备态测量，计数 \((n_1,\ldots,n_4)\) 服从多项分布。求线性无偏估计 \(\hat{\boldsymbol s}\) 的完整协方差矩阵与其迹，讨论纯态、最大混态的数值极限。
(4) 在最大混态 \(\boldsymbol s=0\) 处求关于三个 Bloch 分量的经典Fisher信息矩阵，证明线性估计在该点达到 Cramér–Rao 矩阵界。解释为何靠近物理边界时最大似然估计必须包含正定性约束。
